%%%%%%%%%%%%%%%%%%%%%%%%%%%%%%%%%%%%%%%%%%%%%%%%%%%%%%%%%%%%%%
%%%%%%%%%%%% MA-0404 Herramientas de ciencia de datos I %%%%%%
%%%%%%%%%%%%%%%%%%%%%%%%%%%%%%%%%%%%%%%%%%%%%%%%%%%%%%%%%%%%%%
\newpage
\subsection*{MA-0404 Herramientas de ciencia de datos I}
\addcontentsline{toc}{subsection}{MA-0404 Herramientas de ciencia de datos I}

\hypertarget{MA0404}{}
\begin{refsection}
\begin{table}[H]
\centering{
\begin{tabular}{|ll|}
\hline
\textbf{Año:} II  & \textbf{Requisitos:} MA-0341 Matemática Computacional I \\ 
\textbf{Ciclo:} IV  & \textbf{Correquisitos:} Ninguno \\ 
\textbf{Tipo de curso:} Teórico-práctico & \textbf{Horas semanales:} 5 \\ 
\textbf{Créditos:} 4 & \textbf{Virtualidad:} P, PV \\ \hline
\end{tabular}
}
\end{table}

\subsubsection*{Descripción del curso}

Este curso está dirigido a estudiantes que poseen los conocimientos básicos de programación por lo que se ubica en el segundo año y segundo ciclo. Este curso, junto con Herramientas de Ciencias de Datos II, procura que el/la estudiante aprenda a utilizar y programar herramientas de software que son de uso frecuente en la ciencia de datos, así como en la elaboración de reportes y control de versiones de código. En específico, se busca que el estudiante aprenda a manejar hojas de cálculo, programar en el lenguaje estadístico R, utilizar \LaTeX{} y Markdown e implementar repositorios en línea basados en Git.

\subsubsection*{Objetivo general}

Proporcionar el conocimiento para utilizar y programar herramientas de software para la manipulación y visualización de datos, control de versiones de código y la elaboración de reportes.

\subsubsection*{Objetivos específicos}

Al finalizar el curso el/la estudiante estará en la capacidad de:

\begin{enumerate}
\item Aplicar técnicas de manejo y depuración de datos en hojas de cálculo. (SC04, SC11, SH01.)
\item Usar los fundamentos básicos del lenguaje de programación R en problemas de exploración y visualización de tablas de datos. (SC04, SC11, SH01, SH06.)
\item Utilizar \LaTeX{} y Markdown (RMarkdown) en la elaboración de informes y reportes. (SH09.)
\item Usar comandos básicos de GNU Bash en servidores basados en el sistema operativo Linux. (SH06.)
\item Implementar repositorios en línea basados en el sistema de control de versiones Git. (SH06, SH09.)
\end{enumerate}

\subsubsection*{Contenidos}

\begin{enumerate}
\item \textbf{Manejo y buenas prácticas de datos en hojas de cálculo}
\begin{enumerate}
\item Formateo de tablas de datos en hojas de cálculo.
\item Reducción de errores en el formateo de los datos.
\item Manejo de fechas de hojas de cálculos.
\item Control de calidad y manipulación de datos en hojas de cálculo.
\item Exportar datos desde hojas de cálculo.
\item Estudio de casos sobre las limitaciones de las hojas de cálculo.
\end{enumerate}
\item \textbf{Instalación y configuración de R y RStudio.} Instalación de paquetes en R.
\item \textbf{Introducción al lenguaje de programación R}
\begin{enumerate}
\item Objetos: matrices, listas y \emph{data frames}.
\item Funciones y \emph{scripts}.
\item Bucles y condicionales.
\item Código vectorizado.
\end{enumerate}
\item \textbf{Creación de documentos con \LaTeX{} y RMarkdown.} ProjectTemplate.
\item \textbf{Manejo de datos y visualización con R}
\begin{enumerate}
\item Visualización con \texttt{ggplot2}.
\item Transformación de datos con \texttt{tidyverse}: \texttt{dplyr}, \texttt{tidyr}.
\item Análisis exploratorio de datos: visualización de distribuciones; visualización de asociación entre variables.
\item Manejo de tablas en formato ancho/largo, datos \emph{tidy}.
\item Manejo de cadenas de texto, fechas y factores.
\end{enumerate}
\item \textbf{Introducción a Bash (UNIX).}
\item \textbf{Manejo de Git}
\begin{enumerate}
\item \texttt{git init}, \texttt{git commit}, \texttt{git push}, \texttt{git pull}.
\item Creación de repositorios en línea.
\end{enumerate}
\end{enumerate}

\subsubsection*{Metodología}

Clases magistrales con laboratorios prácticos para demostrar las diferentes herramientas. Además, se le solicitará al estudiante que al finalizar el curso entregue un proyecto basado en un artículo científico o trabajo de investigación de posgrado en donde se le solicitará lo siguiente:

\begin{itemize}
\item Un repositorio \emph{git} (GitHub/GitLab) con la historia de su proyecto.
\item Código de la importación de datos ya sea mediante Excel o R.
\item Código con la transformación de datos.
\item Código con los \emph{scripts} para visualizar sus datos.
\item Reporte descriptivo con los hallazgos encontrados.
\end{itemize}

Incorporar ejemplos de tablas de datos relacionados con información financiera. La persona docente deberá incentivar la lectura de material relacionado con el curso en idioma inglés y el aprendizaje de conceptos básicos de estadística descriptiva a través del descubrimiento de los conceptos en datos reales.

\subsubsection*{Evaluación}

\begin{enumerate}
\item Bitácoras parciales de proyecto.
\item Reporte final del proyecto.
\item Presentación del proyecto.
\end{enumerate}

\subsubsection*{Bibliografía}

\begin{enumerate}
\item Çetinkaya-Rundel, M., \& Hardin, J. (2021). \emph{Introduction to Modern Statistics} (1st ed.). OpenIntro. \url{https://openintro-ims.netlify.app/}
\item Data Organization in Spreadsheets for Ecologists (s/f). \url{https://datacarpentry.org/spreadsheet-ecology-lesson/}
\item Grolemund, G., \& Wickham, H. (2014). \emph{Hands-On Programming with R: Write Your Own Functions and Simulations.} O'Reilly Media. \url{https://rstudio-education.github.io/hopr/}
\item Kelion, L. (2020, 5 de octubre). Excel: Why using Microsoft's tool caused covid-19 results to be lost. BBC News. \url{https://www.bbc.com/news/technology-54423988}
\item Perez-Riverol, Y., et al. (2016). Ten Simple Rules for Taking Advantage of Git and GitHub. \emph{PLOS Computational Biology}, 12(7), e1004947. \url{https://doi.org/10.1371/journal.pcbi.1004947}
\item Wickham, H., \& Grolemund, G. (2016). \emph{R for data science: import, tidy, transform, visualize, and model data} (First edition). O'Reilly.
\item Xie, Y. (2015). \emph{Dynamic Documents with R and knitr} (2nd edition). Chapman and Hall/CRC.
\end{enumerate}

\end{refsection}
